% from ega4-2-fr.tex
\oldpage[IV-2]{109}

\subsection{\texorpdfstring{Modules $Z$-purs et $Z$-clos}{Modules Z-purs et Z-clos}}

Une partie des notions et résultats de cette section et de la suivante sont des
cas particuliers de notions et résultats développés au chapitre III dans la
théorie de la cohomologie locale. Pour la commodité du lecteur, nous en donnons
ici un exposé indépendant.

\begin{env}[5.9.1]
\label{IV.2.s5.9.1-fr}
Soient $X$ un préschéma localement noethérien, $Z$ une partie de $X$
\emph{stable par spécialisation}: cela signifie que pour toute partie finie
$M$ de $Z$, l'adhérence de $M$ est contenue dans $Z$, et par suite $Z$ est
réunion d'une famille filtrante croissante $(Z_\alpha)$ de parties fermées de
$X$; inversement, il est clair qu'une telle réunion est stable par
spécialisation.

Posons $U_\alpha=X-Z_\alpha$, de sorte que $X-Z$ est intersection de la
famille filtrante décroissante d'ouverts $U_\alpha$; soit
$i_\alpha:U_\alpha\to X$ l'injection canonique et pour
$U_\alpha\supset U_\beta$, soit $i_{\alpha\beta}:U_\beta\to U_\alpha$
l'injection canonique, de sorte que l'on a
$i_\beta=i_\alpha\circ i_{\alpha\beta}$. Soit $\mathcal F$ un
$\mathcal O_X$-Module (non nécessairement quasi-cohérent); on a donc
\[
(i_\beta)_*(\mathcal F|U_\beta)
=(i_\alpha)_*((i_{\alpha\beta})_*(\mathcal F|U_\beta));
\]
de l'homomorphisme canonique $(0_{\mathrm I},4.4.3.2)$
\[
\mathcal F|U_\alpha\longrightarrow
(i_{\alpha\beta})_*(\mathcal F|U_\beta)
\]
on déduit donc, par application du foncteur $(i_\alpha)_*$, un homomorphisme
\[
\rho_{\beta\alpha}:(i_\alpha)_*(\mathcal F|U_\alpha)
\longrightarrow(i_\beta)_*(\mathcal F|U_\beta)
\]
et l'on vérifie aussitôt que l'on a
$\rho_{\gamma\alpha}=\rho_{\gamma\beta}\circ\rho_{\beta\alpha}$ pour
$U_\alpha\supset U_\beta\supset U_\gamma$; autrement dit, les
$\mathcal O_X$-Modules $(i_\alpha)_*(\mathcal F|U_\alpha)$ forment un
\emph{système inductif} pour les homomorphismes $\rho_{\beta\alpha}$.
On pose
\begin{equation}
\mathcal H^0_{X/Z}(\mathcal F)
=\varinjlim_\alpha(i_\alpha)_*(\mathcal F|U_\alpha).
\tag{5.9.1.1}\label{IV.2.s5.9.1.1-fr}
\end{equation}

Ce $\mathcal O_X$-Module ne dépend pas de la famille croissante
$(Z_\alpha)$ de fermés dont $Z$ est la réunion: en effet, soit $V$ un ouvert
noethérien de $X$; on sait $(0_{\mathrm{II}},3.10.1)$ que dans la catégorie
des $\mathcal O_V$-Modules, le foncteur
$\mathcal G\mapsto\Gamma(V,\mathcal G)$ commute aux limites inductives; on a
donc en vertu de (5.9.1.1)
\begin{equation}
\Gamma(V,\mathcal H^0_{X/Z}(\mathcal F))
=\varinjlim_\alpha\Gamma(V\cap U_\alpha,\mathcal F).
\tag{5.9.1.2}\label{IV.2.s5.9.1.2-fr}
\end{equation}

Soit alors $(Z'_\lambda)$ une seconde famille filtrante croissante de fermés
de $X$, de réunion $Z$; $V\cap Z_\alpha$ est donc réunion des
$V\cap Z_\alpha\cap Z'_\lambda$; mais $V\cap Z_\alpha$ est localement fermé
dans $X$, donc toute partie fermée irréductible de $V\cap Z_\alpha$ admet un
point générique; comme les $V\cap Z_\alpha\cap Z'_\lambda$ sont fermés dans
$V\cap Z_\alpha$ et forment (pour $\alpha$ fixe) une famille filtrante
croissante, il existe un indice $\lambda$ tel que
$V\cap Z_\alpha\cap Z'_\lambda=V\cap Z_\alpha$
$(0_{\mathrm{III}},9.2.4)$, autrement dit
$V\cap Z_\alpha\subset V\cap Z'_\lambda$. Ceci prouve que les familles
filtrantes décroissantes $V\cap U_\alpha$, $V\cap U'_\lambda$ (où
$U'_\lambda=X-Z'_\lambda$) sont cofinales l'une de l'autre, d'où notre
assertion, en vertu de (5.9.1.2).
\end{env}

